\appendix
\begin{appendices}
\renewcommand{\thefigure}{a\arabic{figure}}
% \crefname{figure}{fig.}{figs.}
% \Crefname{figure}{Fig.}{Figs.}
\setcounter{figure}{0}

\section{Atomic environment}
\label{asec:atomic_environment}
\cref{fig:atom_graph} shows how an atomic environment can be represented by a chemical species-colored and bond length-weighted graph.
\begin{figure}[htbp]
    \centering
    \includegraphics[width=0.3\textwidth]{figs/atomic_graph/atom_graph}
    \caption{A schematic graph representation of the local environment of atom of interest (the middle red atom). The black edges represent the red-red bonds while the red edges represent the red-black bonds. Dashed edges represent the bonds that are not present in the local environment and the atoms outside of the local environment are omitted.}
    \label{fig:atom_graph}
\end{figure}
\newpage

\section{Neural network architecture}
\label{asec:neural_network_architecture}
The structure of the GDyNets framework is depicted in \cref{fig:gdynet}. In the context of this research, the length of the atom feature is designated as $16$. The design employs $6$ graph convolutional layers followed by $6$ dense layers succeeding the pooling layer. The training process spans $30$ epochs with a batch size set at $60$. Optimization is carried out using the default Adam optimizer. For further insights and intricate details, we recommend the reader to refer to the original paper \cite{xie2019graph}.

\begin{figure}[htbp]
    \centering
    \includegraphics[width=0.9\textwidth]{figs/gdynet/gdynet.png}
    \caption{Structure of the GDyNets framework. The two lobes share the same network while reading different data: one reads the trajectory frame at time $t$ and the other reads at time $t+\tau$.}
    \label{fig:gdynet}
\end{figure}
\newpage
    % \begin{figure}[htbp]
    %     \centering
    %     \includegraphics[width=0.75\textwidth]{figs/diffusion/Al_diffusion}
    %     \caption{The variation of the z-positions of Al atoms from the pure Al layer during heat treatment simulation.The colorbar corresponds to the range of the TiAl\textsubscript{3} grain layer, with the lower limit representing the bottom and the upper limit representing the top of the layer. The entire pure Al region lies below the TiAl\textsubscript{3} layer and is shown in dark blue. The "sliced" visualizations in the bottom row present a top-view of the Al atoms with the atoms located above z-position $20$ removed, providing insight into the configuration of Al atoms in the TiAl\textsubscript{3} grain boundaries.}
    %     \label{fig:Al_diffusion}
    % \end{figure}

\section{Distribution of states at 983.5K}
\label{asec:distribution_at_983.5K}
The distributions of the Markov states of Al atoms at \SI{983.5}{\kelvin} are similar to the distributions at \SI{1033.5}{\kelvin}. The results are shown in \cref{fig:Ti_based_983_dot_5K} and \cref{fig:Al_based_983_dot_5K}. The Al$90\%$Ti$10\%$ composition is not included in either system at this temperature due to the slow dynamics and lack of significant dynamic behavior.

\begin{figure}[htbp]
    \def\subplotheight{0.18\textwidth}
    \def\verticalskip{0.5\baselineskip}
    \newcommand{\subplot}[2]{
        \sidesubfloat[]{
            \includegraphics[height=\subplotheight]{#1}\label{#2}
        }
    }

    \subplot{figs/distribution_and_probability/Ti-Al100Ti0-Ti-983_dot_5-4-distribution-A_exam.png}{fig:Ti_distr_sub0_983_dot_5K}
    \subplot{figs/distribution_and_probability/Ti-Al100Ti0-Ti-983_dot_5-4-probability-A_exam.png}{fig:Ti_prob_sub0_983_dot_5K}
    \\
    % [\verticalskip]

    \subplot{figs/distribution_and_probability/Ti-Al95Ti05-Ti-983_dot_5-6-distribution-A_exam.png}{fig:Ti_distr_sub1_983_dot_5K}
    \subplot{figs/distribution_and_probability/Ti-Al95Ti05-Ti-983_dot_5-6-probability-A_exam.png}{fig:Ti_prob_sub1_983_dot_5K}
    
    \caption{The distribution of states of Al atoms in Ti-substrated system at \SI{983.5}{\kelvin}. The substrates are omitted. (a) and (c) depict the spatial distribution of probabilities of Al atoms being in different states for composition Al$100\%$Ti$0\%$ and Al$95\%$Ti$5\%$ respectively. (b) and (d) are the corresponding probability distribution as a function of the z position.}
    \label{fig:Ti_based_983_dot_5K}
\end{figure}

\begin{figure}[htbp]
    \def\subplotheight{0.18\textwidth}
    \def\verticalskip{0.5\baselineskip}
    \newcommand{\subplot}[2]{
        \sidesubfloat[]{
            \includegraphics[height=\subplotheight]{#1}\label{#2}
        }
    }

    \subplot{figs/distribution_and_probability/Al-Al100Ti0-Al-983_dot_5-9-distribution-A_exam.png}{fig:Al_distr_sub0_983_dot_5K}
    \subplot{figs/distribution_and_probability/Al-Al100Ti0-Al-983_dot_5-9-probability-A_exam.png}{fig:Al_prob_sub0_983_dot_5K}
    \\

    \subplot{figs/distribution_and_probability/Al-Al95Ti05-Al-983_dot_5-6-distribution-A_exam.png}{fig:Al-distr_sub1_983_dot_5K}
    \subplot{figs/distribution_and_probability/Al-Al95Ti05-Al-983_dot_5-6-probability-A_exam.png}{fig:Al_prob_sub1_983_dot_5K}
    
    \caption{The distribution of states of Al atoms in Al-substrated system at \SI{983.5}{\kelvin}. The substrates are omitted. (a) and (c) depict the spatial distribution of probabilities of Al atoms being in different states for composition Al$100\%$Ti$0\%$ and Al$95\%$Ti$5\%$ respectively. (b) and (d) are the corresponding probability distribution as a function of the z position.}
    \label{fig:Al_based_983_dot_5K}
\end{figure}
\newpage

\section{Distribution of states and distribution of potential energy}
\label{asec:potential_energy_distribution}
% \tj{This is too brief...}
While it may be tempting to assume that the dynamics are strongly correlated with the atomic potentials, as slow transitions between states could indicate a high energy barrier, this is not always the case, as demonstrated in \cref{fig:atomic_potential}. Specifically, in the Ti-substrated system, there is a single layer of Al atoms in close proximity to the substrate that exhibits significantly lower energy than the atoms located further from the substrate. Conversely, in the Al-substrated system, such a layer does not exist. However, as discussed in \cref{sec:distribution_of_states}, the transition of the solid state Al atoms into other states (and vice versa) is slow in both systems.
\begin{figure}[htbp]
    \centering
    \includegraphics[width=0.5\textwidth]{figs/potential_distribution/potential_distribution}
    \caption{The atomic potential energy of Al-substrated and Ti-substrated systems with $c = \mathrm{Al}95\%\mathrm{Ti}5\%$ at \SI{1033.5}{\kelvin}.}
    \label{fig:atomic_potential}
\end{figure}

\newpage
\section{Transitions between states at 983.5K}
Eigenvectors and timescales for the Al and Ti-substrated systems at \SI{983.5}{\kelvin} are presented in \cref{fig:Al_eigen_983_dot_5K} and \cref{fig:Ti_eigen_983_dot_5K}, respectively. The Al$90\%$Ti$10\%$ composition is not included for the same reason stated in \cref{asec:distribution_at_983.5K}.
\label{asec:timescale_at_983.5K}
\begin{figure}[htbp]
    \def\subplotheight{0.2\textwidth}
    \def\verticalskip{0.5\baselineskip}
    \newcommand{\subplot}[2]{
        \sidesubfloat[]{
            \includegraphics[height=\subplotheight]{#1}\label{#2}
        }
    }
    \subplot{figs/eigen_vectors_and_timescales/Al-Al100Ti0-Al-9-983_dot_5K.png}{fig:Al_eigen_0_983_dot_5K}
    \subplot{figs/eigen_vectors_and_timescales/Al-Al95Ti05-Al-6-983_dot_5K.png}{fig:Al_eigen_1_983_dot_5K}
    
    \caption{The eigen vectors of Al-substrated dynamic systems with (a) Al$100\%$Ti$0\%$, (b) Al$95\%$Ti$5\%$ and (c) Al$90\%$Ti$10\%$ at \SI{983.5}{\kelvin}.}
    \label{fig:Al_eigen_983_dot_5K}
\end{figure}

\begin{figure}[htbp]
    \def\subplotheight{0.2\textwidth}
    \def\verticalskip{0.5\baselineskip}
    \newcommand{\subplot}[2]{
        \sidesubfloat[]{
            \includegraphics[height=\subplotheight]{#1}\label{#2}
        }
    }
    \subplot{figs/eigen_vectors_and_timescales/Ti-Al100Ti0-Ti-4-983_dot_5K.png}{fig:Ti_eigen_0_983_dot_5K}
    \subplot{figs/eigen_vectors_and_timescales/Ti-Al95Ti05-Ti-6-983_dot_5K.png}{fig:Ti_eigen_1_983_dot_5K}
    
    \caption{The eigen vectors of Ti-substrated dynamic systems with (a) Al$100\%$Ti$0\%$, (b) Al$95\%$Ti$5\%$ and (c) Al$90\%$Ti$10\%$ at \SI{983.5}{\kelvin}.}
    \label{fig:Ti_eigen_983_dot_5K}
\end{figure}

\end{appendices}
